\chapter{Optimizaci\'on de Consultas, Planes de Ejecuci\'on e \'indices}
\label{cap:optimizacion_explain_indices}

\section{Marco Conceptual y Te\'orico (Curr\'iculo Universitario)}

El optimizador de consultas basado en costos (\textit{Cost-Based Optimizer / CBO}) de PostgreSQL genera m\'ultiples planes de ejecuci\'on y selecciona aquel con el menor costo estimado en unidades arbitrarias de acceso a disco y CPU.

\subsection{Componentes de un Plan de Ejecuci\'on}
\begin{itemize}
    \item \textbf{Sequential Scan (\texttt{Seq Scan}):} Lectura completa de todas las p\'aginas de la tabla. Eficiente para tablas peque\~nas o cuando la consulta recupera un alto porcentaje del total de filas ($>20\%$).
    \item \textbf{Index Scan:} B\'usqueda en la estructura B-Tree del \'indice seguida de la lectura de las p\'aginas de datos en el Heap.
    \item \textbf{Index Only Scan:} Recuperaci\'on de datos directamente desde las hojas del \'indice sin acceder al Heap (requiere que el Visibility Map confirme que la p\'agina no tiene tuplas modificadas pendientes).
    \item \textbf{Bitmap Index Scan + Bitmap Heap Scan:} Combina m\'ultiples \'indices generando un mapa de bits en memoria antes de leer el Heap secuencialmente por bloques.
\end{itemize}

\subsection{Estrategias de Indexaci\'on Avanzada}
\begin{enumerate}
    \item \textbf{\'indices Compuestos:} Creados sobre m\'ultiples columnas \texttt{(col1, col2)}. Siguen la regla del prefijo izquierdo.
    \item \textbf{\'indices Parciales:} \texttt{CREATE INDEX ... WHERE condicion;} optimizan espacio e indexan \'unicamente filas relevantes (ej. \texttt{WHERE status = 'PENDIENTE'}).
    \item \textbf{\'indices Expresionales:} \texttt{CREATE INDEX ... ON users(LOWER(email));}.
    \item \textbf{\'indices Cubrientes (\texttt{INCLUDE}):} A\~naden columnas al nivel de hoja sin incluirlas en la clave de ordenamiento del B-Tree.
\end{enumerate}

---

\section{Casos de Estudio y Pr\'actica Aplicada}

\begin{lstlisting}[style=sqlstyle, caption={An\'alisis comparativo de costos con EXPLAIN ANALYZE}]
-- 1. Inspecci\'on de consulta sin \'indice adecuado
EXPLAIN (ANALYZE, BUFFERS, TIMING, FORMAT TEXT)
SELECT order_id, customer_id, order_date
FROM orders
WHERE customer_id = 'ALFKI'
  AND order_date >= '1997-01-01';

-- 2. Creaci\'on de \'indice compuesto cubriente
CREATE INDEX idx_orders_customer_date_inc
ON orders (customer_id, order_date)
INCLUDE (order_id);

-- 3. Re-evaluaci\'on del plan de ejecuci\'on (Transformaci\'on a Index Only Scan)
EXPLAIN (ANALYZE, BUFFERS, TIMING, FORMAT TEXT)
SELECT order_id, customer_id, order_date
FROM orders
WHERE customer_id = 'ALFKI'
  AND order_date >= '1997-01-01';
\end{lstlisting}

---

\section{Taller de Consolidaci\'on del Cap\'itulo}

\begin{businessproblem}
Una consulta cr\'itica del panel administrativo de ventas experimenta degradaci\'on de rendimiento severa al escanear secuencialmente millones de registros en la tabla de detalle de \'ordenes. Se requiere identificar cuellos de botella mediante \texttt{EXPLAIN}, proponer la estrategia de indexaci\'on y verificar la reducci\'on del costo computacional.
\end{businessproblem}

\subsection{Soluci\'on T\'ecnica y Estrategia de Tuning}

\begin{lstlisting}[style=sqlstyle, caption={Optimizaci\'on de consulta anal\'itica y creaci\'on de \'indice GIN / Parcial}]
-- Consulta objetivo: B\'usqueda de \'ordenes de alto volumen con descuento mayor al 15%
EXPLAIN (ANALYZE, BUFFERS)
SELECT
    order_id,
    product_id,
    unit_price,
    quantity,
    discount
FROM order_details
WHERE discount > 0.15
ORDER BY quantity DESC;

-- Estrategia: \'indice parcial con ordenamiento invertido en cantidad
CREATE INDEX idx_order_details_discount_qty
ON order_details (quantity DESC)
WHERE discount > 0.15;
\end{lstlisting}

\begin{engtip}
\begin{itemize}
    \item \textbf{Impacto de Buffers (\texttt{shared hit}):} La m\'etrica \texttt{Buffers: shared hit=N read=M} indica cu\'antos bloques de 8KB fueron le\'idos directamente desde la memoria RAM (\textit{Shared Buffers}) versus disco f\'isico. Una optimizaci\'on exitosa maximiza el ratio de aciertos (\texttt{hit percent} $>99\%$).
\end{itemize}
\end{engtip}
